\begin{problem}[Uniqueness of sheafification from the universal property]
Suppose $(\mathcal S,\eta)$ and $(\mathcal T,\theta)$ are two sheafifications of the same presheaf $\mathcal F$. Prove there is a unique isomorphism $\alpha:\mathcal S\to\mathcal T$ satisfying $\alpha\eta=\theta$.
\end{problem}

\paragraph{Why this problem matters.}
A universal construction is useful only if its output is canonical up to unique compatible isomorphism.

\paragraph{Useful ideas before starting.}
Use the universal property, stalk preservation, and the fact that a morphism of sheaves is an isomorphism when it is so on every stalk.  Keep the distinction between a presheaf statement on sections and a local statement on germs explicit.

\paragraph{Guided hints.}
Factor $\theta$ through $\eta$ and $\eta$ through $\theta$, then compare composites with identities.

\ifdefined\FullProblemDossiers
\begin{solution}
Because $\mathcal T$ is a sheaf, the universal property of $(\mathcal S,\eta)$ gives a unique $\alpha:\mathcal S\to\mathcal T$ with $\alpha\eta=\theta$. Similarly there is a unique $\beta:\mathcal T\to\mathcal S$ with $\beta\theta=\eta$. Then
\[
(\beta\alpha)\eta=\beta\theta=\eta.
\]
Both $\beta\alpha$ and $\id_{\mathcal S}$ are maps $\mathcal S\to\mathcal S$ whose composition with $\eta$ is $\eta$; uniqueness forces $\beta\alpha=\id$. Similarly $\alpha\beta=\id$. Hence $\alpha$ is the required unique isomorphism.
\end{solution}

\paragraph{What happened mathematically.}
No comparison of explicit constructions is needed: universality itself supplies and identifies the isomorphism.

\paragraph{Alternative approaches.}
Whenever possible one may replace the espace-etale model by the two-plus construction, or conversely replace a cover/refinement argument by the universal property.  The two approaches agree because both construct the same associated sheaf.

\paragraph{Extensions.}
Apply this result to compare the plus and espace-etale constructions.

\paragraph{Related problems.}
Compare with VI.14.P1 for the complete legacy construction and with the later sheaf chapters for operations that are deliberately not migrated here.

\paragraph{Provenance.}
New synthesis problem distilled from the uniqueness clause in the legacy construction.
\else
\paragraph{Provenance.}
New synthesis problem distilled from the uniqueness clause in the legacy construction.
\fi
